% =============================================================================
\section{Tarefa 1 --- Modelos dos sensores}
\label{sec:modelos}
% =============================================================================

Designa-se por $\{s\}$ o referencial do sensor e por $\{w\}$ um referencial
fixo, com $z$ vertical para cima; $R$ é a orientação de $\{s\}$ em $\{w\}$.
As grandezas medidas levam til.

\paragraph{Acelerómetro.} Não mede a aceleração $\vect a$ do sensor, mas a
força específica: a aceleração menos a gravidade
$\vect g = (0;\,0;\,-g)$, vista no referencial do sensor. Parado, mede $1\,g$
para cima. O modelo adotado é
\begin{equation}
\med a = S_a\,R^{T}\,(\vect a - \vect g) + \vect b_a + \vect n_a ,
\label{eq:acc}
\end{equation}
com $\vect b_a$ o desvio (\emph{bias}), $S_a$ uma matriz $3\times3$ com os
fatores de escala na diagonal e o desalinhamento entre eixos fora dela, e
$\vect n_a$ ruído branco de variância $\sigma_a^2$ por eixo.

\paragraph{Giroscópio.} Mede a velocidade angular $\vect\omega$ do sensor,
expressa em $\{s\}$:
\begin{equation}
\med\omega = S_g\,\vect\omega + \vect b_g + \vect n_g ,
\label{eq:gyro}
\end{equation}
com os mesmos três tipos de erro. O desvio $\vect b_g$ não é constante: muda
cada vez que o sensor é ligado e com a temperatura~\cite{titterton2004}.

\paragraph{Montagem.} O integrado não está alinhado com o suporte. A rotação
$R_{ms}$ entre os eixos do sensor e os da ferramenta do robô, e a posição
$\vect r$ do sensor em relação ao eixo de rotação, entram nas medições: um
acelerómetro a uma distância $\vect r$ do eixo mede a mais
$\vect\alpha\times\vect r + \vect\omega\times(\vect\omega\times\vect r)$.

\paragraph{Parâmetros a estimar.} Tabela~\ref{tab:param}.

\begin{table}[H]
\centering
\caption{Parâmetros dos modelos e ensaio de onde se tira cada um.}
\label{tab:param}
\small
\begin{tabularx}{\textwidth}{@{}lLcl@{}}
\toprule
\rowcolor{tabcab}
Parâmetro & Significado & N.º & Ensaio \\
\midrule
$\vect b_a$ & desvio do acelerómetro & 3 & B, B2 \\
$S_a$ & escala (e desalinhamento) do acelerómetro & 3 (+3) & B, B2 \\
$\sigma_a$ & ruído do acelerómetro & 3 & A \\
$\vect b_g$ & desvio do giroscópio & 3 & A, e início de cada ensaio \\
$S_g$ & escala do giroscópio & 3 & C \\
$\sigma_g$ & ruído do giroscópio & 3 & A \\
$R_{ms}$, $\vect r$ & orientação e posição do sensor no suporte & 3 + 3 & B2 e C, no robô \\
\bottomrule
\end{tabularx}
\end{table}

\begin{destaque}
\textbf{Porque é que estes parâmetros importam:} as Tarefas 4 e 6 integram as
medições no tempo. Um desvio $b_g$ no giroscópio dá um erro de orientação
$b_g\,t$; um desvio $b$ na aceleração dá um erro de posição
$\tfrac12\,b\,t^2$. Erros pequenos e constantes deixam de ser pequenos.
\end{destaque}
